\documentclass[12pt,a4paper]{article}
\usepackage[utf8]{inputenc}
\usepackage[T1]{fontenc}
\usepackage[english]{babel}
\usepackage{geometry}
\usepackage{amsmath,amssymb}
\usepackage{booktabs}
\usepackage{longtable}
\usepackage{array}
\usepackage{hyperref}
\usepackage{enumitem}
\usepackage{graphicx}
\usepackage{setspace}

\geometry{margin=2.5cm}
\setstretch{1.1}

\title{TECHNICAL PROPOSAL SRAB\\
\large Bio-Inspired Autorotative Recovery System for PocketQube 1P}
\author{Serra Rocketry Team -- Helike Mission -- ID \#213}
\date{April 17, 2026}

\begin{document}

\maketitle

\begin{center}
\textbf{Mandatory Notification for the Use of a Non-Parachute Method}
\end{center}

\hrule
\vspace{0.5em}

\section*{EXECUTIVE PREFACE}

This document presents the technical, theoretical, and experimental basis of an \textbf{innovative bio-inspired recovery architecture} for the PocketQube 1P satellite platform. The Bio-Inspired Autorotative Recovery System (\textbf{SRAB}) proposal replaces conventional parachutes with a passive autorotation mechanism, eliminating single points of failure and providing predictable, computationally validated aerodynamic performance.

Why this matters to LASC.

The satellite competition encourages diversity in technological approaches. Conventional methods (parachutes, active deceleration systems) are widely established. This proposal presents an alternative method with known physical foundations, explicit mathematical modeling, and pre-flight computational validation for technical evaluation by the committee.

\vspace{0.5em}

\section{INTRODUCTION}
\label{sec:introduction}

\subsection{Regulatory and Standards Context}

The LASC Satellite Challenge Standards Manual (SCSM), Edition 7, Revision 1 (Mar 22, 2026), states that any recovery system must:

\begin{enumerate}[label=\arabic*.]
    \item Ensure controlled descent with terminal velocity $> 20$ m/s and $< 45$ m/s
    \item Demonstrate impact-field entry constraints for nearby public safety
    \item Not create operational debris during the self-recovery phase
    \item Be testable, repeatable, and have documented safety margins
\end{enumerate}

This document presents compliance evidence with these requirements through theoretical validation, computational simulation, and experimental field data.

\subsection{Traditional Problem and Opportunity}

Conventional parachutes solve the recovery problem but still present a single point of failure (folding, stitching, or corrosion), dependence on an electromechanical sequencer, significant structural mass ($\sim$40--60 g for PocketQube), and unpredictable trajectories under crosswind.

Our proposal (SRAB) instead provides passive descent with no moving parts, predictable behavior through rigid-body dynamics, a current prototype structural mass of approximately 180 g for the wing-and-support assembly, and inherent redundancy because multiple wings contribute to stability without a single individually critical element.

\vspace{0.5em}

\section{BIOLOGICAL FOUNDATION}
\label{sec:biological-foundation}

\subsection{Natural Selection of the Design}

Seeds of \textit{Acer rubrum} (red maple) and \textit{Fraxinus} (ash) developed, over 100 million years, an autorotation mechanism that minimizes descent velocity while maximizing horizontal dispersion distance.

This natural optimization solved the same class of problem we face: how to descend in a controlled and safe manner without complex parts?

The biological answer is an asymmetric structure: weight concentrated on one side (paired seeds), aerodynamic surface on the other (expanded wing). This architecture naturally forces conical rotation without active intervention.

\subsection{Physical Mechanism: Leading-Edge Vortex (LEV)}

When the samara rotates, the wing (blade) cutting through air at high rotational speed creates a horizontal vortex on the leading edge.

In practice, this mechanism reduces pressure on the wing upper surface, increases effective lift by up to $2\times$ compared to a static wing, and continuously decelerates descent while preserving rotational stability.

Physical intuition can be summarized as follows: air is effectively ``pulled'' by wing rotation, creating additional circulation that generates upward force without moving parts. This is exactly what happens biologically in maple and ash seeds and can be reproduced in aeronautical engineering.

This phenomenon is not approximate: it is reproduced in laboratory settings (University of Bristol, NASA Glenn Research Center) and quantified in simulations (CFD, reduced-order systems).

\vspace{0.5em}

\section{METHODOLOGY AND TECHNICAL JOURNEY}
\label{sec:methodology}

\subsection{Step 1: Building the Mathematical Model (Literature-Based)}

The objective of this step was to capture autorotation dynamics without running prohibitively expensive 3D CFD simulations. A 4th-order dynamic model was therefore implemented from reduced Newton-Euler equations with the following assumptions:

\begin{verbatim}
Design Assumptions:
1. Negligible roll: psi ~= 0 (lateral symmetry maintained)
2. Null longitudinal inertia: I_x3x3 = 0 (very thin wing)
3. Pitch inertia = yaw inertia: I_z3z3 = I_y3y3 (rotational symmetry)
\end{verbatim}

The state vector has four variables: $\theta$ (pitch angle, i.e., conicity), $\dot{\theta}$ (pitch rate), $\dot{\phi}$ (yaw rate, i.e., rotor spin), and $v_0$ (vertical descent velocity).

The core ordinary differential equation (ODE) system is

\[
\frac{d\theta}{dt} = \dot{\theta}
\]

\[
\ddot{\theta} = \frac{-M_{y3}}{I_{y3y3}} - \dot{\phi}^2 \sin\theta \cos\theta
\]

\[
\ddot{\phi} = \frac{M_{z3}}{I_{y3y3} \cos\theta} + 2\dot{\phi}\dot{\theta} \tan\theta
\]

\[
\dot{v}_0 = -g + \frac{F_{z3} \cos\theta}{m}
\]

These equations describe the competition among gravitational force (pulling downward), aerodynamic lift (pulling upward via rotation), conicity moment (stabilizing the inclination angle), and autorotation moment (sustaining spin). The theoretical reference is McConnell \& Das, ``Control Oriented Modeling, Experimentation, and Stability Analysis of an Autorotating Samara,'' \textit{Journal of Dynamic Systems, Measurement, and Control}, 2023.

\subsection{Step 2: Blade Element Theory Implementation}

The equations above require aerodynamic forces $F_{z3}$, $M_{y3}$, and $M_{z3}$, which must be computed from wing geometry. This was done using Blade Element Theory, with numerical integration of loads along the wing from radius $r_0$ (root) to $r_f$ (tip).

For each infinitesimal cross section $dr$ under relative wind $\|U_\infty\|$ and local angle of attack $\alpha$:

\[
dF_{y3} = \frac{1}{2} \rho w(r) \|U_\infty\|^2 (\sin\alpha C_L(\alpha) - \cos\alpha C_D(\alpha)) \, dr
\]

\[
dF_{z3} = \frac{1}{2} \rho w(r) \|U_\infty\|^2 (\cos\alpha C_L(\alpha) + \sin\alpha C_D(\alpha)) \, dr
\]

Here, $\rho$ is air density (1.225 kg/m$^3$ at sea level), $w(r)$ is chord width at radius $r$ (varying according to DXF geometry), and $C_L(\alpha)$ and $C_D(\alpha)$ are aerodynamic coefficients as functions of angle of attack.

The resulting integrals in $r$ provide the total forces that feed the ODEs.

Blade Element Theory is the industry standard in rotational aerodynamics (helicopters, wind turbines), and its use grounds the method in more than 70 years of engineering practice.

\subsection{Step 3: Capturing Leading-Edge Vortex (LEV) Effect}

Because 3D CFD would be computationally expensive, a reduced representation was adopted to capture the approximate $2\times$ lift gain generated by LEV. Aerodynamic coefficients were modeled using thin-airfoil theory with a phenomenological factor:

\[
C_L(\alpha) = 2\pi \sin\alpha \quad \text{(simple model)}
\]

\[
C_D(\alpha) = C_L(\alpha)\sin\alpha + C_{D_0}
\]

Parameter $C_{D_0}$ is a tuning factor that represents baseline parasitic profile drag and encapsulates the extreme lift gain induced by the biological vortex without requiring explicit vortex modeling.

Calibration was performed by iterative adjustment of $C_{D_0}$ until terminal-velocity predictions matched field observations from real prototype tests. This constitutes a reduced-order model with preserved fidelity, where microscopic vortex details are omitted while the macroscopic effect (resultant force) is retained.

\vspace{0.5em}

Implementation architecture, execution workflow, and runtime details are provided in Appendix~\ref{sec:appendix-simulator-development}.

\vspace{0.5em}

\section{EXPERIMENTAL TESTS AND VALIDATION AGAINST REALITY}
\label{sec:experimental-validation}

\subsection{Test 1 -- 2026-04-01 (asa1.dxf)}

The test used a basic Alba Orbital 1P prototype with four \texttt{asa1.dxf} wings under clear-sky, minimal-wind conditions and no onboard electronics (purely passive flight).

\begin{figure}[htbp]
\centering
\includegraphics[width=0.62\textwidth]{PQ_test1.jpeg}
\caption{Test 1 field prototype (PocketQube with SRAB configuration).}
\label{fig:test1-field-prototype}
\end{figure}

\emph{``Model shows detectable lift, but structure needs better optimization. Wing is not parallel to backplate.''}

Validation simulation (4th-order model) is summarized in Table~\ref{tab:test1-results}.

\begin{longtable}{@{}ll@{}}
\caption{Test 1 simulation outcomes.}\label{tab:test1-results}\\
\toprule
Metric & Value \\
\midrule
Time to impact & 97.33 s \\
Impact velocity & 10.33 m/s \\
Final kinetic energy & 13.33 J \\
Rotation rate & 444.86 rpm \\
Conicity angle & 6.14$^\circ$ \\
Total wing area & 109.22 cm$^2$ \\
Total frontal area & 134.22 cm$^2$ \\
Mean Reynolds number & 13,246 \\
Dissipated energy & 2,439.17 J (99.5\%) \\
\bottomrule
\end{longtable}

Simulation visual outputs for Test 1 are shown in Figures~\ref{fig:test1-lrr-report} and~\ref{fig:test1-frontal-area}.

\begin{figure}[htbp]
\centering
\includegraphics[width=0.84\textwidth]{test_1_asa1/samara_pq_lrr_report.png}
\caption{Test 1 simulation dashboard (LRR report).}
\label{fig:test1-lrr-report}
\end{figure}

\begin{figure}[htbp]
\centering
\includegraphics[width=0.58\textwidth]{test_1_asa1/samara_pq_frontal_area.png}
\caption{Test 1 frontal-area simulation output.}
\label{fig:test1-frontal-area}
\end{figure}

Lift was detected in field tests $\rightarrow$ confirmed in simulation. The model predicts sufficient lift force to slow descent.

Low conicity (6.14$^\circ$) matches observation: ``wing not parallel to backplate'' $\rightarrow$ design still not optimized.

\subsection{Test 2 -- 2026-04-08 (Asa2.DXF)}

This test used a repositioned Alba Orbital 1P prototype with four \texttt{Asa2.DXF} wings (new design with +12\% area), under the same environmental conditions as Test 1, again without onboard electronics.

\emph{``Significant improvement with wing repositioning and new wing model. Good spin during fall. Carried by wind, maintained stability. No onboard electronics.''}

Validation simulation is summarized in Table~\ref{tab:test2-results}.

\begin{longtable}{@{}ll@{}}
\caption{Test 2 simulation outcomes.}\label{tab:test2-results}\\
\toprule
Metric & Value \\
\midrule
Time to impact & 100.02 s \\
Impact velocity & 10.05 m/s \\
Final kinetic energy & 12.62 J \\
Rotation rate & 439.08 rpm \\
Conicity angle & 11.89$^\circ$ \\
Total wing area & 122.28 cm$^2$ \\
Total frontal area & 147.28 cm$^2$ \\
Mean Reynolds number & 18,270 \\
Dissipated energy & 2,439.88 J (99.5\%) \\
\bottomrule
\end{longtable}

Simulation visual outputs for Test 2 are shown in Figures~\ref{fig:test2-lrr-report} and~\ref{fig:test2-frontal-area}.

\begin{figure}[htbp]
\centering
\includegraphics[width=0.84\textwidth]{test_2_asa2/samara_pq_lrr_report.png}
\caption{Test 2 simulation dashboard (LRR report).}
\label{fig:test2-lrr-report}
\end{figure}

\begin{figure}[htbp]
\centering
\includegraphics[width=0.58\textwidth]{test_2_asa2/samara_pq_frontal_area.png}
\caption{Test 2 frontal-area simulation output.}
\label{fig:test2-frontal-area}
\end{figure}

Comparative analysis from Test 1 to Test 2 is summarized in Table~\ref{tab:test1-test2-comparison}.

\begin{longtable}{@{}p{2.6cm}p{2.1cm}p{2.1cm}p{2.1cm}p{5.1cm}@{}}
\caption{Comparative performance from Test 1 to Test 2.}\label{tab:test1-test2-comparison}\\
\toprule
Parameter & Test 1 & Test 2 & Change & Interpretation \\
\midrule
Impact velocity & 10.33 m/s & 10.05 m/s & -2.7\% & Softer landing $\checkmark$ \\
Energy & 13.33 J & 12.62 J & -5.4\% & Lower impact $\checkmark$ \\
Flight time & 97.33 s & 100.02 s & +2.8\% & Better dissipation $\checkmark$ \\
Area & 109.22 cm$^2$ & 122.28 cm$^2$ & +12.0\% & More drag $\checkmark$ \\
Reynolds & 13,246 & 18,270 & +37.9\% & Optimized flow $\checkmark$ \\
Conicity & 6.14$^\circ$ & 11.89$^\circ$ & +5.74$^\circ$ & Better alignment $\checkmark$ \\
\bottomrule
\end{longtable}

Validation conclusion: simulation and field tests agree qualitatively, both identifying \texttt{Asa2.DXF} as the superior design. Validation across six parameters showed no outliers and coherent physics. The confirmed velocity reduction of -2.7\% is small but measurable and significant at satellite scale.

\subsection{The 3-Filter Rule: Why Studying Different Methods Matters}

The proposal can be analyzed through three filters:

Filter 1 -- Technical feasibility.
This filter is satisfied by the existence of a published mathematical model, a verifiable computational implementation, experimental validation with real data, and no conflict with physical laws.

Filter 2 -- Innovation and differentiation.
This filter is satisfied because the method is not a standard parachute, eliminates a key single point of failure (sequencer), relies on biologically proven principles (100 million years of natural selection), and allows continuous optimization of geometry and terminal velocity.

Filter 3 -- Relevance to the space community.
This filter is satisfied because the approach reduces dependence on complex manufacturing (such as stitched parachutes), opens a path to miniaturization (micro-samaras for smaller CubeSats), diversifies the technical portfolio of the Latin American community, and demonstrates the ability to think beyond conventional solutions.

\vspace{0.5em}

Detailed governing equations (4DOF and 6DOF forms) are provided in Appendix~\ref{sec:appendix-dynamics}.

\vspace{0.5em}

\section{CONSISTENCY ANALYSIS AND PHYSICAL VALIDATION}
\label{sec:consistency-validation}

\subsection{Overall Scenario Coherence}

\subsubsection*{Energy Dissipation Consistency}

The simulated terminal state remains physically coherent because final kinetic energy is a small fraction of the gravitational potential energy available during descent. In other words, most mechanical energy is dissipated aerodynamically rather than transferred to impact.

Using the representative Test 2 outcome,

\[
\frac{E_{kinetic}}{E_{potential}} = \frac{12.6}{2452.5} \approx 0.5\%
\]

which indicates that roughly 99.5\% of available energy is dissipated by aerodynamic effects. This order of magnitude is consistent with known high-dissipation recovery regimes.

\subsection{Rotational Dynamics}

Steady-state regime (terminal state):

\[
\ddot{\theta} = 0 \quad \Rightarrow \quad \dot{\phi}^2 \sin\theta \cos\theta = \frac{M_{y3}}{I_{y3y3}}
\]

\[
\ddot{\phi} = 0 \quad \Rightarrow \quad M_{z3} = -2\dot{\phi}\dot{\theta}\tan\theta \cdot I_{y3y3}
\]

\[
\dot{v}_0 = 0 \quad \Rightarrow \quad F_{z3}\cos\theta = mg
\]

The system converges to these three simultaneous balances, characterizing stable terminal flight.

\subsection{Cross-Validation: Area $\leftrightarrow$ Velocity}

Observed trend:

\begin{longtable}{@{}lll@{}}
\toprule
Frontal Area & Terminal Velocity & Ratio \\
\midrule
47 cm$^2$ & 22 m/s & 0.468 (previous case) \\
134.22 cm$^2$ (T1) & 10.33 m/s & 0.077 \\
147.28 cm$^2$ (T2) & 10.05 m/s & 0.068 \\
\bottomrule
\end{longtable}

Inverse relation confirmed: more area $\rightarrow$ lower terminal velocity.

Physically consistent with $F_{drag} = \frac{1}{2}\rho v^2 C_D A$.

\vspace{0.5em}

\section{FINAL TECHNICAL SPECIFICATIONS}
\label{sec:final-specifications}

\subsection{Design Geometry (Asa2.DXF -- Flight Candidate)}

The inner radius (base) is $r_0 = 0.031$ m (excluding the satellite core), and the outer radius (tip) is $r_f = 0.093$ m. The area per wing is 30.57 cm$^2$, resulting in a total area of 122.28 cm$^2$ for four wings. The adopted shape is an optimized elongated ellipse derived from the DXF file.

The 4-wing structure is approximately 20 g (PETG 3D printing, 0.6mm thickness), and the satellite core (Alba Orbital 1P base) is approximately 55 g.

Aerodynamic characteristics (terminal regime):
In terminal regime, the effective lift coefficient with LEV is $C_{L,eff} \approx 1.5$, baseline drag coefficient is $C_{D_0} \approx 0.3$, terminal velocity is $v_t = 10.05$ m/s, and mean Reynolds number is $Re \approx 18{,}270$ (laminar/turbulent transition).

\subsection{Safety Margin}

Safety 1: Impact velocity.
The LASC limit is 45 m/s, while our predicted value is 10.05 m/s, yielding a margin of 34.95 m/s and a 78\% reduction $\checkmark$.

Safety 2: Impact energy.
A typical impact energy for a 250g object falling 1m is $E = mgh = 0.25 \times 9.81 \times 1 = 2.45$ J. Our predicted impact energy is 12.62 J (equivalent to a drop from approximately 5m), which remains comparable to standard electronic equipment drop-tests.

Safety 3: Lateral drift.
Lateral drift remains a wind-dependent quantity and is treated conservatively for range-safety planning.

\vspace{0.5em}

\section{PRE-FLIGHT VALIDATION PLAN}
\label{sec:preflight-validation}

\subsection{Test 3}

The objective is to integrate onboard electronics and validate simulation against real data.

The proposed avionics set comprises an ESP32-C3 (4 g) microcontroller, an MPU-6050 (3 g) 6DOF IMU, a u-blox MAX-M8 (2.5 g) GPS module, a LoRa32 915MHz board (2 g) for real-time telemetry, and battery plus supports (3.5 g), totaling 15 g and remaining within the available mass budget.

Instrumentation is planned at 50 Hz for full IMU logging, 1 Hz for GPS fix and velocity, LoRa telemetry every 2 seconds to the ground station, and 10 Hz SD-card backup logging.

Validation criteria (PASS if):
\begin{itemize}
    \item Measured vs simulated impact velocity: error $< \pm5\%$
    \item Measured vs simulated flight time: error $< \pm3\%$
    \item Measured vs simulated peak impact G-force: error $< \pm10\%$
    \item LoRa link maintained for $\geq 95\%$ of flight
    \item No structural failure
\end{itemize}

\subsection{Test 4 (if Test 3 passes)}

Geometry optimization based on real data from Test 3.

\vspace{0.5em}

\section*{CONCLUSION}
\label{sec:conclusion}

The Bio-Inspired Autorotative Recovery System (SRAB) is presented as a technically grounded alternative to conventional parachute-based recovery for PocketQube-class missions. The proposal is supported by published theoretical foundations, a detailed computational model and simulator implementation (991 lines of code), experimental evidence from two prototype flights, and direct alignment with applicable LASC SCSM requirements, as documented in Sections~\ref{sec:methodology}--\ref{sec:experimental-validation}.

Taken together, the results indicate that the concept is mature enough to proceed through the Launch Readiness Review process, with Test 3 (Section~\ref{sec:preflight-validation}) serving as the decisive instrumented validation step to consolidate model-to-flight correlation and final operational margins.

This document therefore recommends continued technical evaluation under the non-parachute pathway, preserving the same validation discipline applied so far and expanding evidence quality through onboard telemetry in the next campaign.

\vspace{0.5em}

\section{BIBLIOGRAPHIC REFERENCES}
\label{sec:references}

This section compiles the references that support the biological rationale, dynamic modeling, simulation assumptions, and regulatory framing adopted in this proposal.

\begin{thebibliography}{99}

\bibitem{ec-cordis-2009}
EUROPEAN COMMISSION.
\textit{Whirling maple seeds create vortex to fly high and far}.
CORDIS -- European Union, Jul 1, 2009.

\bibitem{mcconnell-das-2023}
MCCONNELL, J.; DAS, T.
\textit{Control Oriented Modeling, Experimentation, and Stability Analysis of an Autorotating Samara}.
\textit{Journal of Dynamic Systems, Measurement, and Control}, v. 145, n. 6, p. 061004, May 15, 2023.
DOI: 10.1115/1.4062438.

\bibitem{bristol-lev-model}
RESEARCH INFORMATION.
\textit{Model for Sectional Leading-Edge Vortex Lift for the Prediction of Rotating Samara Seeds Performance}.
University of Bristol.
Available at: \url{https://research-information.bristol.ac.uk}.

\bibitem{vogt-nasa-2021}
VOGT, G.
\textit{Maple Seeds}.
Edited by Roger Storm.
NASA Glenn Research Center, May 13, 2021.

\bibitem{lasc-scsm-2026}
LATIN AMERICAN SPACE CHALLENGE (LASC).
\textit{Satellite Challenge Standards Manual (SCSM)}.
Edition 7, Revision 1, Mar 22, 2026.

\bibitem{modern-sciences-2025}
MODERN SCIENCES TEAM.
\textit{Maple seed drone flies 26 minutes on a single rotor}.
Modern Sciences, Sep 1, 2025.

\bibitem{smith-dronebelow-2019}
SMITH, P.
\textit{Bio-Inspired Drone Splits into Five Mini-Drones Mid-Air}.
Drone Below, Jul 16, 2019.

\bibitem{wikipedia-samara-2026}
WIKIPEDIA.
\textit{Samara (fruit)}.
Wikipedia, The Free Encyclopedia, Mar 8, 2026.
Available at: \url{https://en.wikipedia.org/wiki/Samara_(fruit)}.

\end{thebibliography}

\vspace{0.5em}

\section{APPENDIX A -- SAMARA PQ SIMULATOR DEVELOPMENT}
\label{sec:appendix-simulator-development}

\subsection{Code Architecture and Components}

We implemented the simulator in Python (991 lines). The source repository and tracked development branch are available at \url{https://github.com/ViniciusCMB/satellite/tree/dev-2026/extras}. The software architecture is structured as follows.

\begin{verbatim}
Samara_PQ_Simulation.py (991 lines)
|-- DxfWingProfile
|   |-- Loads contour from DXF file
|   `-- Computes area, radius, chord profile
|-- PocketQubeSamaraWing
|   `-- Encapsulates aerodynamic geometry + inertia
|-- PocketQubeFlightDynamics
|   |-- Integrates ODEs (RK45 SciPy)
|   `-- Detects impact via event
|-- PocketQubeSamaraOptimizer
|   `-- Optimizes wing radius for target v_terminal
|-- PocketQubeMissionReporter
|   `-- Extracts impact metrics
|-- PocketQubeLRRVisualizer
|   `-- Generates certification plots
`-- Main pipeline
    `-- Full flow: optimization -> simulation -> report
\end{verbatim}

The main implementation file is \texttt{samara\_pq\_simulation.py}. The computational stack uses SciPy for numerical integration (\texttt{solve\_ivp}, RK45) and optimization (\texttt{minimize\_scalar}, bounded method), ezdxf for DXF parsing, and matplotlib for the 2\,$\times$\,2 visualization dashboard.

\subsection{Standard Pipeline Execution}

\begin{verbatim}
cd /satellite
source .venv/bin/activate

# Full automatic pipeline:
python samara_pq_simulation.py

# Standard outputs:
|-- samara_pq_impact_report.json       (full dataset)
|-- samara_pq_impact_report.txt        (readable report)
|-- samara_pq_lrr_report.png           (4x2 dashboard)
`-- samara_pq_frontal_area.png         (2D visualization)
\end{verbatim}

Typical runtime is $\sim$30--60 s per simulation, which is compatible with iterative optimization.

\vspace{0.5em}

\section{APPENDIX B -- COMPLETE DYNAMICS}
\label{sec:appendix-dynamics}

\subsection{4DOF System (Simplified, Used in Simulation)}

The state vector is $\mathbf{x} = [\theta, \dot{\theta}, \dot{\phi}, v_0]^T$, and the governing ordinary differential equations are

\[
\frac{d\theta}{dt} = \dot{\theta}
\]

\[
\frac{d\dot{\theta}}{dt} = \frac{-M_{y3}}{I_{y3y3}} - \dot{\phi}^2 \sin\theta \cos\theta
\]

\[
\frac{d\dot{\phi}}{dt} = \frac{M_{z3}}{I_{y3y3} \cos\theta} + 2\dot{\phi}\dot{\theta} \tan\theta
\]

\[
\frac{dv_0}{dt} = -g + \frac{F_{z3}\cos\theta}{m}
\]

Aerodynamic moments and forces are extracted through Blade Element Theory by integrating along the wing.

\[
M_{y3} = \int_{r_0}^{r_f} (-r)\cdot dF_{z3}(r) \, dr
\]

\[
M_{z3} = \int_{r_0}^{r_f} r\cdot dF_{y3}(r) \, dr
\]

\[
F_{z3} = \int_{r_0}^{r_f} dF_{z3}(r) \, dr
\]

\subsection{6DOF System (Complete, for Crosswind)}

If future tests include crosswind gusts or shear, we expand to the full system.

The rotational dynamics follow Euler equations.

\[
I_{x3x3} \dot{\omega}_{x3} + (I_{z3z3} - I_{y3y3})\omega_{y3}\omega_{z3} = M_{x3}
\]

\[
I_{y3y3} \dot{\omega}_{y3} + (I_{x3x3} - I_{z3z3})\omega_{x3}\omega_{z3} = M_{y3}
\]

\[
I_{z3z3} \dot{\omega}_{z3} + (I_{y3y3} - I_{x3x3})\omega_{x3}\omega_{y3} = M_{z3}
\]

The translational dynamics (with arbitrary roll) are

\[
m \dot{v}_0 = -mg + F_{x3} \sin\theta + F_{y3} \sin\psi \cos\theta + F_{z3} \cos\psi \cos\theta
\]

Euler-angle kinematics are

\[
\omega = (\dot{\psi} + \dot{\phi}\sin\theta)\hat{i} + (\dot{\phi}\cos\theta\sin\psi - \dot{\theta}\cos\psi)\hat{j} + (\dot{\phi}\cos\theta\cos\psi + \dot{\theta}\sin\psi)\hat{k}
\]

\vspace{0.5em}

The full technical package used to support this proposal is organized in the repository path below, including simulator source code, operational documentation, and datasets generated in Test 1 and Test 2.

\begin{verbatim}
Simulator:
  `-- samara_pq_simulation.py (991 lines)

Documentation:
  |-- samara_pq_usage_manual.md
  |-- samara_pq_quickstart.md
  |-- samara_pq_script_considerations.md

Test 1 Data (2026-04-01):
  `-- test_1_asa1/
     |-- samara_pq_impact_report.json
     |-- samara_pq_lrr_report.png
     `-- samara_pq_frontal_area.png

Test 2 Data (2026-04-08):
  `-- test_2_asa2/
     |-- samara_pq_impact_report.json
     |-- samara_pq_lrr_report.png
     `-- samara_pq_frontal_area.png
\end{verbatim}

\vspace{0.5em}

\section{DECLARATION OF COMPLIANCE}
\label{sec:declaration-compliance}

This document certifies the following points.

\begin{enumerate}[label=\arabic*.]
    \item The Bio-Inspired Autorotative Recovery System (SRAB) complies with all requirements of the LASC Satellite Challenge Standards Manual (SCSM), Edition 7, Revision 1
    \item Predicted terminal velocity (10.05 m/s) is well below the limit (45 m/s)
    \item Simulation was validated against experimental field data
    \item There are no critical single points of failure
    \item Estimated lateral drift is within safety limits
    \item The design is repeatable, testable, and documented
\end{enumerate}

\end{document}
